\section{Conclusion and Future Work}
\label{sec:conclusion}

We introduced a novel self-supervised framework for in-the-wild video reshooting that addresses the severe scarcity of paired multi-view data. By generating pseudo multi-view training triplets solely from monocular videos, our approach bypasses the need for explicit 3D annotations. We achieve this through a token-concatenation architecture within a DiT-based model, utilizing Offset RoPE and targeted augmentations to effectively route dynamic source textures. Extensive evaluations demonstrate that our method achieves state-of-the-art temporal consistency and geometric control, confirming that the model natively learns implicit 4D structures directly from 2D data. 

While our approach significantly advances dynamic reshooting, we identify key areas for future exploration. First, concatenating source tokens doubles the sequence length, impacting generation speed. Because source video information remains static across diffusion timesteps, future work could explore sophisticated KV-caching mechanisms to largely eliminate this computational overhead. Second, extreme trajectories that move entirely outside the original scene's boundaries yield blank anchor videos, removing geometric guidance. To resolve this, future iterations could employ a hybrid conditioning scheme leveraging our synthetic dataset mixture to co-train the model on both visual anchors and explicit camera poses. Finally, to further mitigate blank anchors and even more intricate camera paths, an autoregressive framework could be explored where the anchor condition is generated by forward-warping the previously generated target frame. This would naturally improve conditioning stability, representing a critical step toward true 4D generative world models.